\chapter{Fundraising, Dilution \& Runway}\label{ch:fundraising-and-dilution}

% Scope: cap tables & dilution (pre/post-money), SAFEs & convertible notes, term sheets &
%   liquidation preferences, investor metrics & the ARR multiple. NEW chapter (closes the
%   founder-finance gap). Follows valuation (\cref{ch:valuation}) and capital structure
%   (\cref{ch:risk-and-capital-structure}): from enterprise value to how founders actually finance.
% Cross-links: eq:scenario-ev (ch.20 bear-case anchor for Seed pre-money); eq:dcf / eq:ev-revenue
%   (valuation chain reused as investor frame); eq:rule40 / eq:burn-multiple / eq:runway (ch.21/26
%   efficiency metrics that drive the fundraising narrative).

\section{Cap Tables \& Ownership Dilution}\label{sec:cap-table}
% complexity: 3

What does raising a round actually cost the founders? The cap table --- the ledger recording
every share, option, and convertible security --- answers that question precisely, and its
arithmetic is unforgiving. Understanding dilution before signing means knowing the
post-money formula cold.

When a company issues new shares to investors at an agreed valuation, the existing owners'
percentage shrinks even as the total value of their stake may rise. The two quantities a
founder must track are (i) the \emph{post-money valuation}, which fixes the price per share,
and (ii) the \emph{founder ownership percentage} after the new shares are absorbed into the
fully diluted count. The core identity is:
\begin{equation}\label{eq:post-money}
\begin{gathered}
  \mathrm{PostMoney} \;=\; \mathrm{PreMoney} \;+\; \mathrm{Investment}, \\
  \text{Investor\%} \;=\; \frac{\mathrm{Investment}}{\mathrm{PostMoney}}, \qquad
  \text{Founder\%}_{\mathrm{post}} \;=\; \text{Founder\%}_{\mathrm{pre}}\,\bigl(1 - \text{Investor\%}\bigr).
\end{gathered}
\end{equation}
Founders begin at \qty{100}{\percent} and dilute multiplicatively with every round. Each
successive financing multiplies the surviving percentage by \((1 - \text{new investor\%})\);
the percentages of all existing holders --- founders, angels, and prior-round investors ---
shrink proportionally to make room for the incoming capital.

A subtlety that \cref{eq:post-money} conceals is the \emph{option pool shuffle}. Investors
typically require that the unallocated employee option pool be expanded \emph{before} the new
money is counted, which means that dilution falls entirely on founders rather than being
shared with the new investor \parencite{feld2019venture}. A term sheet requiring a
\qty{15}{\percent} post-money pool from a \money{10000000} pre-money company effectively
forces the founder to absorb the pool dilution before the price per share is struck. Founders
who miss this arithmetic pay for it twice.

\begin{example}[Two-round cap table]
The company's scenario-probability-weighted enterprise value was estimated at approximately
\money{16500000} (see \cref{eq:scenario-ev}), but the bear case of \money{8000000} is the
conservative anchor for a Seed pre-money negotiation. The company raises a Seed round at a
\money{6000000} pre-money valuation on \money{2000000} of investment.

Applying \cref{eq:post-money}:
\[
  \mathrm{PostMoney}_{\text{Seed}} \;=\; \money{6000000} + \money{2000000} \;=\; \money{8000000}.
\]
The Seed investor's stake: \(\money{2000000}/\money{8000000} = \qty{25}{\percent}\). The
founder's stake falls from \qty{100}{\percent} to \qty{75}{\percent}.

Now the company raises a Series~A at a \money{10000000} pre-money valuation, issuing
\money{2000000} of new equity.
\[
  \mathrm{PostMoney}_{\text{Series A}} \;=\; \money{10000000} + \money{2000000} \;=\; \money{12000000}.
\]
The Series~A investor's stake: \(\money{2000000}/\money{12000000} \approx \qty{16.7}{\percent}\).
Applying the multiplicative rule from \cref{eq:post-money}:
\[
  \text{Founder\%}_{\text{post-A}}
  \;=\; \qty{75}{\percent} \;\times\; (1 - 0.167)
  \;=\; \qty{75}{\percent} \;\times\; 0.833
  \;\approx\; \qty{62.5}{\percent}.
\]
The cumulative post-Series~A cap table: founder \qty{62.5}{\percent}; Seed investor \qty{20.8}{\percent};
Series~A investor \qty{16.7}{\percent}. Total: \qty{100}{\percent}.

The implication is not that the founder ``lost'' \qty{37.5}{\percent}: on a \money{12000000}
post-money, their \qty{62.5}{\percent} is worth \money{7500000} --- more than the entire
company was valued at pre-Seed.
\end{example}

\begin{admonition}{Warning}
The percentage a founder holds after a round is almost irrelevant in isolation. What matters
is the absolute outcome: \qty{62.5}{\percent} of a \money{200000000} exit is worth more than
\qty{90}{\percent} of a \money{10000000} exit by a factor of 11. Optimising for ownership
retention at the cost of capital that accelerates growth is a common founder error. The
correct frame is the expected value of the founder's dollar position, not the fraction of the
pie --- a distinction the DCF logic of \cref{eq:dcf} makes rigorous.
\end{admonition}

Dilution compounds across rounds as the product of survival fractions; option pools, warrants,
and convertible instruments further expand the denominator. The full picture requires tracking
the \emph{fully diluted} share count rather than issued shares alone.
\textcite{metrick2021venture} develop the complete multi-round mechanics, including option-pool
and warrant adjustments, and show that late-stage founders with strategic investors often
retain between \qty{10}{\percent} and \qty{30}{\percent} at IPO.

\section{SAFEs \& Convertible Notes}\label{sec:safes-convertibles}
% complexity: 3

Early-stage companies rarely have the operating history to justify a hard pre-money
negotiation. Two instruments defer that conversation: the Simple Agreement for Future Equity
(SAFE), introduced by Y Combinator in 2013, and the older convertible note. Both inject cash
immediately and convert into equity at a future priced round, but their economic structures
differ materially. A convertible note is debt --- it accrues interest, carries a maturity
date, and can force a founder into default if no priced round arrives in time. A SAFE is not
debt; it carries no interest and no maturity. The economic protection for early investors
comes instead from two parameters: a \emph{valuation cap} (the maximum valuation at which
the SAFE converts, regardless of how high the Series~A is priced) and a \emph{discount rate}
(a percentage reduction off the Series~A price per share). Most SAFEs use one or both.

The conversion price --- the price per share the SAFE investor pays --- is:
\begin{equation}\label{eq:safe-conversion}
  P_{\mathrm{conv}}
  \;=\; \min\!\bigl(\mathrm{Cap},\;\mathrm{SeriesA}_{\mathrm{pre}}\bigr)
        \;\times\; (1 - \delta),
  \qquad
  N_{\mathrm{new}}
  \;=\; \frac{\mathrm{Investment}}{P_{\mathrm{conv}}},
\end{equation}
where \(\delta\) is the discount rate and \(N_{\mathrm{new}}\) is the number of new shares
issued to the SAFE holder upon conversion. When the Series~A pre-money exceeds the cap, the
cap binds; when it falls below the cap, the Series~A price (discounted by \(\delta\)) governs.
The bound that produces the larger share count --- and therefore the lower effective price per
share --- will always be the one the investor prefers.

\begin{example}[SAFE conversion with cap and discount]
An angel invests \money{500000} on a SAFE with a \money{5000000} cap and a \qty{20}{\percent}
discount. The company subsequently raises a Series~A at a \money{12000000} pre-money
valuation.

First, determine which bound governs:
\begin{itemize}
  \item Cap-based effective valuation: \money{5000000}.
  \item Discount-based effective valuation: \(\money{12000000} \times (1 - 0.20) = \money{9600000}\).
\end{itemize}
The cap (\money{5000000}) is lower, so the cap binds --- the SAFE investor converts as if
the pre-money were only \money{5000000}. Applying \cref{eq:safe-conversion} with
\(\delta = 0\) once the cap governs (the cap already embeds the economic benefit; using the
cap and the discount simultaneously would double-count):
\[
  P_{\mathrm{conv}} \;=\; \money{5000000} \;\times\; 1.00
                   \;=\; \money{5000000} \text{ (effective valuation).}
\]
The SAFE holder's ownership fraction at conversion:
\[
  \text{SAFE\%}
  \;=\; \frac{\mathrm{Investment}}{\mathrm{Cap}}
  \;=\; \frac{\money{500000}}{\money{5000000}}
  \;=\; \qty{10}{\percent} \;\text{(on a pre-Series-A, post-SAFE basis).}
\]
After the Series~A investor takes \qty{16.7}{\percent} (as in the prior example), all existing
holders --- SAFE investor plus founders --- dilute proportionally:
\[
  \text{SAFE\%}_{\text{post-A}}
  \;=\; \qty{10}{\percent} \;\times\; (1 - 0.167)
  \;\approx\; \qty{8.3}{\percent}.
\]
The SAFE investor ends up at \qty{8.3}{\percent} of a \money{12000000} post-money, an
implied value of roughly \money{1000000} on a \money{500000} investment --- a \qty{2}{\times}
mark-up at the moment of the Series~A close, before any exit.
\end{example}

\begin{admonition}{Warning}
Multiple SAFEs at different caps and different discounts convert at different prices,
issuing different numbers of shares. Founders who raise a \money{250000} SAFE at a
\money{3000000} cap, then a \money{500000} SAFE at a \money{5000000} cap, and then a
\money{1000000} SAFE at an \money{8000000} cap must model the full conversion waterfall
before signing the Series~A term sheet --- the interaction can be far more dilutive than
any individual SAFE implied. Each SAFE is solved independently using \cref{eq:safe-conversion},
and the resulting share counts stack before the new institutional money is admitted.
\end{admonition}

SAFEs and notes trade cheapness and speed against complexity at the conversion event. The
economics are governed entirely by the cap and discount parameters; the legal structure is
standardised \parencite{feld2019venture}. Once the Series~A is priced, the SAFE disappears
into the cap table as preferred stock, subject to the same liquidation mechanics described in
the next section.

\section{Term Sheets \& Investor Economics}\label{sec:term-sheets}
% complexity: 3

A venture term sheet allocates cash-flow rights, control rights, and downside protection
simultaneously. The economic core is the \emph{liquidation preference}: the right of
preferred stockholders to recover their investment (or a multiple of it) before common
shareholders receive anything in a sale, merger, or dissolution. Understanding the preference
is understanding who gets paid in every realistic exit scenario.

Two preference structures dominate. In \emph{non-participating preferred}, the investor
chooses the better of two outcomes: take the stated preference or convert to common stock and
share pro-rata. In \emph{participating preferred}, the investor takes the preference
\emph{and} then participates pro-rata in the remainder --- the ``double-dip.'' Non-participating
is the market standard; participating appears in competitive markets and distress rounds, and
is materially worse for founders. The formula for exit proceeds to the investor is:
\begin{equation}\label{eq:liquidation-preference}
  \mathrm{Investor}_{\text{non-part}}
  \;=\; \max\!\bigl(\mathrm{Pref},\; p \cdot \mathrm{Exit}\bigr),
  \qquad
  \mathrm{Investor}_{\text{part}}
  \;=\; \mathrm{Pref} \;+\; p \cdot (\mathrm{Exit} - \mathrm{Pref}),
\end{equation}
where \(\mathrm{Pref} = m \times \mathrm{Investment}\) is the liquidation preference
(multiple \(m\), typically \(1\times\)) and \(p\) is the investor's ownership fraction.
Founders receive what remains after investors are satisfied.

\begin{example}[Non-participating vs.\ participating preference]
A Series~A investor puts in \money{5000000} for \qty{25}{\percent} of the company. The
liquidation preference is \(1\times\) non-participating. The company sells for \money{8000000}.

\emph{Non-participating}: the investor's preference is \money{5000000}; pro-rata conversion
would yield \(25\% \times \money{8000000} = \money{2000000}\). Applying \cref{eq:liquidation-preference}:
\[
  \mathrm{Investor}_{\text{non-part}}
  \;=\; \max(\money{5000000},\;\money{2000000})
  \;=\; \money{5000000}.
\]
Founders and common stockholders receive \(\money{8000000} - \money{5000000} = \money{3000000}\).

\emph{Participating}: the investor first takes \money{5000000}, then participates in the
remainder:
\[
  \begin{aligned}
    \mathrm{Investor}_{\text{part}}
    &\;=\; \money{5000000} \;+\; 0.25 \times (\money{8000000} - \money{5000000}) \\
    &\;=\; \money{5000000} \;+\; \money{750000} \;=\; \money{5750000}.
  \end{aligned}
\]
Founders receive \(\money{8000000} - \money{5750000} = \money{2250000}\). The participation
clause transferred \money{750000} from founders to the investor --- the ``double-dip'' is
\money{750000} in this scenario, or roughly \qty{25}{\percent} of the incremental upside
above the preference threshold.

The magnitude of the double-dip grows with exit size: at a \money{20000000} exit the
participating investor would take \(\money{5000000} + 25\%\times\money{15000000} =
\money{8750000}\), versus \(\max(\money{5000000},\,\money{5000000}) = \money{5000000}\)
non-participating. Participation always costs founders at intermediate exits.
\end{example}

\begin{admonition}{Warning}
Anti-dilution provisions compound the preference structure. Broad-based weighted-average
anti-dilution --- the market standard --- adjusts the conversion price of existing preferred
shares proportionally if a down round occurs, limiting, but not eliminating, the effective
repricing. Full-ratchet anti-dilution reprices the \emph{entire} prior preferred block to the
new low price regardless of volume, which can obliterate founder equity in a single down
round. The punitive arithmetic makes full-ratchet provisions a founder-unfriendly red flag;
insist on broad-based weighted-average \parencite{feld2019venture}. In distressed markets,
the combination of a \(1\times\) participating preference and full-ratchet anti-dilution can
leave common stockholders with near-zero proceeds even at moderate exit values.
\end{admonition}

Kaplan and Strömberg's empirical analysis of venture contracts shows that the allocation of
cash-flow rights (preferences, anti-dilution), control rights (board seats, voting vetoes),
and liquidation rights are deliberately decoupled to align incentives under
uncertainty \parencite{kaplan2003financial}. A preference that looks protective at round
close can become punitive at exit; modelling the full waterfall across plausible exit
scenarios before signing is not optional.

\section{Investor Metrics \& the Fundraising Narrative}\label{sec:investor-metrics}
% complexity: 4

How should a founder frame the company's value to an investor? The answer is not the
narrative alone; it is the narrative \emph{grounded in the numbers} that investors actually
use to benchmark and price startups. In 2026 the primary quantitative filter is not growth
rate in isolation but the relationship between growth and capital consumption, expressed
through the Rule of 40 (\cref{eq:rule40}) and the burn multiple (\cref{eq:burn-multiple}).
The valuation anchor, in turn, is the EV/ARR multiple from \cref{eq:ev-revenue}:
\begin{equation}\label{eq:arr-multiple}
  \mathrm{PreMoney} \;\approx\; \mathrm{ARR} \;\times\; \frac{\mathrm{EV}}{\mathrm{ARR}},
\end{equation}
where the EV/ARR multiple (\cref{eq:ev-revenue}) is the market's current clearing price per
dollar of recurring revenue, adjusted for the company's growth-efficiency profile. A
high-growth, capital-efficient company commands a premium multiple; a slow-growing, high-burn
company faces a steep discount. The Rule of 40 score and the burn multiple are the two
levers investors use to move the multiple \parencite{metrick2021venture}.

\begin{figure}[htb]
  \centering
  \includegraphics[width=0.92\linewidth]{cap-table-waterfall}
  \caption{%
    \textbf{Cap-table dilution and liquidation waterfall.}
    \emph{Left panel:} founder ownership across three financing rounds (Pre-seed
    \qty{100}{\percent} $\to$ Seed \qty{75}{\percent} $\to$ Series~A \qty{62.5}{\percent}).
    \emph{Right panel:} exit proceeds to founders vs.\ the Series~A investor under
    non-participating and participating \(1\times\) liquidation preferences, across exit
    values from \money{5000000} to \money{30000000} (Series~A: \money{5000000} invested,
    \qty{25}{\percent} ownership). The participation kink is visible around \money{20000000},
    where the investor's conversion value overtakes the preference.%
  }
  \label{fig:cap-table-waterfall}
\end{figure}

\begin{example}[Investor-metrics narrative for a Series~A pitch]
The running company: ARR = \money{4000000}; trailing ARR growth = \qty{30}{\percent};
gross margin = \qty{70}{\percent}. In isolation, \qty{30}{\percent} growth at this ARR
could support an \(8\times\) EV/ARR multiple, implying a pre-money of
\[
  \money{4000000} \;\times\; 8 \;=\; \money{32000000}.
\]
But the investor looks at \cref{eq:rule40}: Rule of 40 score = \qty{30}{\percent} growth
$+ (-\qty{15}{\percent})$ EBITDA margin $= 15$. A score of 15 is well below the 40-point
threshold where growth and profitability are balanced; the investor discounts the multiple
from $8\times$ to $6\times$:
\[
  \mathrm{PreMoney} \;=\; \money{4000000} \;\times\; 6 \;=\; \money{24000000}.
\]
The burn multiple (\cref{eq:burn-multiple}): the company added \money{1200000} net new ARR
(30\% of \money{4000000}) while burning \money{1500000} net cash, yielding a burn multiple
of \(\money{1500000}/\money{1200000} = 1.25\times\). At $1.25\times$ the burn multiple is
acceptable --- below the $2\times$ threshold that triggers material haircuts --- and the
founder should lead with it. Combined with the runway forecast from \cref{eq:runway} (cash
/ net burn rate), the narrative becomes: ``We have 18 months of runway at current burn, a
$1.25\times$ burn multiple, and a clear path to Rule of 40 parity by month 24 as gross
margin leverage kicks in.''

The \money{24000000} pre-money sits comfortably between the base and bull DCF scenarios of
\cref{eq:scenario-ev} (\money{15000000}--\money{28000000}), providing investors with a
market-comparable anchor for the DCF they will run independently. Cross-linking the
multiple to the DCF removes the appearance of arbitrary pricing and aligns the conversation
around fundamentals.
\end{example}

\begin{admonition}{Note}
ARR multiples vary significantly by sector, growth tier, and market conditions; the
\(6\times\) figure in the example reflects a moderate-efficiency, mid-growth profile in the
2026 environment. Best-in-class companies scoring above 60 on the Rule of 40 have commanded
multiples of \(10\times\)--\(15\times\) ARR. The appropriate multiple is not a formula
output but a market observation; the EV/ARR frame of \cref{eq:ev-revenue} provides the
theoretical anchor (growth rate, margin, and discount rate determine value), while comparable
transactions provide the empirical check.
\end{admonition}

The fundraising narrative is a translation layer: it renders the DCF logic of \cref{eq:dcf}
into the operating metrics investors actually track. Rule of 40, burn multiple, and ARR
multiple are not alternatives to discounted-cash-flow thinking; they are shortcuts to the
same present-value conclusion. A founder who understands why \cref{eq:arr-multiple} is just
a compressed version of \cref{eq:ev-revenue} --- which is itself a compressed version of
\cref{eq:dcf} --- can defend any metric a sophisticated investor challenges, because the
chain leads back to first principles. \textcite{metrick2021venture} document how institutional
investors build their own DCF models alongside the multiples, using the gap between the two
as a signal of whether the founder understands their own business.
